% IEEE-style LaTeX draft for an approximately 8-page two-column framework paper.
% If IEEEtran.cls is installed, this source uses IEEEtran conference mode.
% If IEEEtran.cls is not installed, it falls back to a two-column article class
% so the draft can still be compiled locally for editing and length checks.
% If the target special call is a journal special issue, change [conference]
% to [journal] in the IEEEtran branch and adjust the author block.
\newif\ifieeetran
\IfFileExists{IEEEtran.cls}{%
  \ieeetrantrue
  \documentclass[conference]{IEEEtran}
}{%
  \ieeetranfalse
  \documentclass[10pt,twocolumn]{article}
}

\usepackage{amsmath,amssymb}
\ifieeetran\else
\setlength{\textwidth}{7.0in}
\setlength{\textheight}{9.25in}
\setlength{\oddsidemargin}{-0.25in}
\setlength{\evensidemargin}{-0.25in}
\setlength{\topmargin}{-0.45in}
\setlength{\columnsep}{0.22in}
\fi

\hyphenation{Simulink micro-grid super-visory}

\begin{document}

\title{A Modular Framework for Guarded Autonomous Digital Twin Supervisory Control}

\ifieeetran
\author{
\IEEEauthorblockN{Author Name(s) Omitted for Draft}
\IEEEauthorblockA{Department / Laboratory / Institution\\
City, Country\\
email@example.com}
}
\else
\author{Author Name(s) Omitted for Draft\\
Department / Laboratory / Institution, City, Country\\
\texttt{email@example.com}}
\date{}
\fi

\maketitle

\begin{abstract}
Digital twins are increasingly used to represent, monitor, and evaluate
engineered systems, yet many deployments remain advisory rather than
autonomous. Moving from passive digital twin visualization toward autonomous
digital twin control requires more than connecting an artificial intelligence
model to a simulator. A practical autonomous digital twin controller must
define a stable state-action contract, separate high-level supervisory
decisions from low-level plant control, enforce safety constraints outside the
policy model, support multiple runtime backends, and produce an auditable
record of each decision. This paper proposes a modular framework for guarded
autonomous digital twin supervisory control. The framework treats the digital
twin as an environment exposed through a compact observation interface and a
bounded command interface. A supervisory runner coordinates observation,
policy selection, command limiting, backend execution, safety checking, and
mission reporting. Control policies are interchangeable and may include
deterministic rule-based policies, learned policies, or large-language-model
decision policies, while an independent guard layer enforces hard constraints,
soft-constraint reporting, command clipping, rate limiting, and fallback
behavior. The framework is demonstrated through a battery-assisted dc load-bus
digital twin with MATLAB/Simulink and UDP runtime interfaces. The case study
illustrates how the proposed architecture supports mock testing, live digital
twin communication, policy substitution, and transparent decision logging. The
result is a reusable design pattern for developing autonomous digital twin
controllers that are inspectable, extensible, and constrained by explicit
operational safety contracts.
\end{abstract}

\ifieeetran
\begin{IEEEkeywords}
Digital twin, autonomous control, supervisory control, guarded autonomy,
large language models, Simulink, microgrids, dc power systems, safety
constraints, decision logging.
\end{IEEEkeywords}
\else
\noindent\textbf{Keywords---}
Digital twin, autonomous control, supervisory control, guarded autonomy,
large language models, Simulink, microgrids, dc power systems, safety
constraints, decision logging.
\par\medskip
\fi

\section{Introduction}

Digital twins have become an important mechanism for representing physical
systems, integrating simulation models with operational data, and supporting
engineering decision-making. In power and energy systems, manufacturing,
transportation, and cyber-physical infrastructure, digital twins can provide a
computational view of asset state, forecast future behavior, and evaluate
hypothetical interventions before those interventions are applied to a real
system. Despite this promise, many digital twin systems remain passive: they
visualize state, answer questions, or provide offline analysis, but they do not
participate directly in closed-loop or supervisory control.

Autonomous digital twin control is a natural next step, but it also creates a
systems problem that is different from model construction alone. A digital
twin controller must interact with a model or live simulation, interpret state
estimates, select actions, and apply those actions through a communication
pathway. If artificial intelligence (AI) is used as part of the decision
process, the architecture must also define what the AI is allowed to decide,
how its outputs are validated, and how the system responds when an AI output
is unsafe, malformed, delayed, or unavailable. In practical engineering
environments, the controller must also remain auditable: operators and
researchers should be able to inspect what the digital twin observed, what
command was issued, why the command was selected, which constraints were
active, and why the mission succeeded or failed.

These requirements suggest that autonomous digital twin control should not be
treated as a direct connection between an AI model and a plant model. Instead,
it should be treated as a guarded supervisory-control framework. The digital
twin provides a structured observation. A policy proposes a bounded high-level
command. A guard layer validates and limits the command. A backend adapter
applies the command to a simulation, live model, or physical interface. A
reporting layer records the complete decision trace. This architecture allows
autonomy to be introduced incrementally while preserving clear responsibility
boundaries between the policy, the digital twin, the safety logic, and the
plant-level controllers.

This paper proposes such a framework. The framework is designed for autonomous
digital twin supervisory control, where the autonomous agent is responsible
for high-level operating decisions rather than direct low-level control
signals. In the motivating case study, a digital twin of a battery-assisted dc
load bus exposes state variables such as state of charge (SOC), bus voltage,
battery power, load power, source power, source availability, and fault flags.
The supervisory controller chooses from a bounded command set: charge, hold,
discharge, or safe. The inner converter control remains inside the Simulink
model, while the autonomous supervisor commands a battery-power reference.
This separation prevents the autonomous policy from issuing raw actuator
commands such as duty-ratio values.

The proposed framework is intentionally policy agnostic. A deterministic
state-of-charge and bus-voltage policy can be used for baseline behavior. A
learned or large-language-model (LLM) policy can be substituted through the
same interface. In all cases, the runner owns hard safety stops, command
clipping, rate limits, and mission logging. Thus, the framework supports
experimentation with advanced policy models without transferring safety
responsibility to those models.

The main contributions of this paper are:
\begin{itemize}
    \item a modular framework for autonomous digital twin supervisory control
    based on an explicit observation, goal, command, and backend contract;
    \item a guarded execution architecture that separates policy proposals
    from safety enforcement, command limiting, and mission termination;
    \item a policy interface that supports deterministic, learned, or
    LLM-based supervisory decisions without changing the digital twin backend;
    \item a backend abstraction that supports mock simulation,
    MATLAB/Simulink engine execution, and live UDP communication with a
    running digital twin; and
    \item a case-study implementation for a battery-assisted dc bus digital
    twin, including decision traces and mission-report generation for
    framework evaluation.
\end{itemize}

\section{Motivation and Requirements}

\subsection{Digital Twins as Runtime Control Contexts}

A digital twin is commonly understood as a digital representation of a
physical system that is linked to data, models, or operational context
\cite{dt_foundation,dt_survey}. In many engineering applications, the digital
twin supports monitoring, diagnosis, simulation, and forecasting. Once the
digital twin is connected to a control task, however, its role changes. It is
no longer only an information source; it becomes part of a decision loop.

This shift introduces several architectural questions. What information does
the twin expose to the controller? Which actions may the controller request?
How are constraints represented? Does the digital twin own the plant-level
controller, or does the autonomous agent command low-level actuator variables
directly? How are decisions logged? How does the system behave when the
communication path fails?

These questions are not solved by model fidelity alone. A high-fidelity
digital twin can still be unsafe or difficult to control if its interface to
the autonomous supervisor is unstructured. Conversely, a simplified digital
twin can be useful for early autonomy research if it exposes a clear
state-action contract and supports repeatable testing. The proposed framework
therefore focuses on the interface layer between digital twin runtime and
autonomous supervisory policy.

\subsection{Supervisory Control Rather than Direct Actuation}

The framework proposed here is designed for supervisory control. The
autonomous digital twin controller selects high-level commands, operating
modes, or references, while lower-level control loops remain responsible for
fast plant dynamics. For a battery-connected dc bus, the supervisor may
request a battery-power reference, but it should not directly command the
converter duty ratio. The duty ratio, pulse-width modulation, current
limiting, and converter dynamics remain inside the plant controller or
simulation model.

This separation is particularly important when AI models are involved. A
language model or learned policy may be useful for selecting a mode,
interpreting a fault context, or choosing a conservative power request. It
should not be responsible for raw switching commands. The framework therefore
constrains the policy action space to bounded supervisory commands and
delegates fast physical control to the digital twin or embedded control layer.

\subsection{Framework Requirements}

The framework is organized around seven requirements:
\begin{enumerate}
    \item \emph{Explicit state-action contract}: the twin exposes typed
    observations and accepts bounded commands.
    \item \emph{Policy substitutability}: deterministic, learned, or LLM
    policies can be exchanged without changing the runner or backend.
    \item \emph{Safety ownership outside the policy}: hard constraints,
    command bounds, rate limits, and fallback behavior are enforced by the
    runner rather than by the policy alone.
    \item \emph{Backend independence}: the same supervisory loop can operate
    over mock simulation, step-based Simulink execution, or live UDP
    communication.
    \item \emph{Communication robustness}: stale observations, repeated
    commands, packet schemas, startup transients, and backend exceptions are
    handled by adapters and guards.
    \item \emph{Auditability}: each decision records the observation,
    selected command, command reason, status, and active violations.
    \item \emph{Incremental autonomy}: development can progress from
    deterministic policies and mock simulations toward AI policies and live
    digital twin communication.
\end{enumerate}

These requirements position the contribution as a framework paper rather than
as a claim that a particular controller is globally optimal. The objective is
to define and demonstrate an architecture that makes autonomous digital twin
supervisory control modular, inspectable, and safe enough for systematic
experimentation.

\section{Proposed Framework}

\subsection{Architecture Overview}

Fig.~\ref{fig:architecture} summarizes the proposed architecture. The
framework consists of six components: a digital twin backend, a structured
observation object, a goal object, a policy module, a guarded runner, and a
reporting layer. At each supervisor tick, the runner receives or requests an
observation from the backend. It evaluates hard and soft constraints. If a
hard constraint is violated, the runner sends a safe command and terminates
the mission. If the mission has satisfied its success criteria for the
required duration, the runner returns success. Otherwise, it calls the active
policy, limits the resulting command, applies the command through the backend,
records the event, and advances to the next observation.

\begin{figure*}[t]
\centering
\fbox{\begin{minipage}{0.94\textwidth}
\centering
\vspace{0.5em}
Digital Twin Backend
$\rightarrow$
Observation Contract
$\rightarrow$
Guarded Runner
$\rightarrow$
Policy Interface
$\rightarrow$
Command Limiter
$\rightarrow$
Backend Command Path
$\rightarrow$
Mission Log
\vspace{0.5em}

\smallskip
\small
Backends: mock simulator, MATLAB/Simulink engine bridge, live UDP interface.
Policies: deterministic rule policy, learned policy, LLM policy.
Guards: hard stops, soft-violation logging, command clipping, rate limits, fallback behavior.
\vspace{0.5em}
\end{minipage}}
\caption{Proposed guarded autonomous digital twin supervisory-control framework.
The policy proposes high-level commands, while the runner enforces safety and
execution constraints before commands reach the digital twin backend.}
\label{fig:architecture}
\end{figure*}

\subsection{Formal Representation}

Let the digital twin at supervisory step $k$ expose an observation
\begin{equation}
    o_k = O(DT, t_k),
\end{equation}
where $O$ is the backend observation function, $DT$ is the digital twin
runtime, and $t_k$ is the current simulation or wall-clock time. Let the
mission goal be
\begin{equation}
    G = \{B_{\mathrm{soft}}, B_{\mathrm{hard}}, T_{\mathrm{success}}, T_{\max}\},
\end{equation}
where $B_{\mathrm{soft}}$ defines desired operating bands,
$B_{\mathrm{hard}}$ defines hard safety limits, $T_{\mathrm{success}}$
defines the required stable hold time, and $T_{\max}$ defines the maximum
mission duration.

A policy proposes a command:
\begin{equation}
    u_k^p = \pi(o_k, G).
\end{equation}
The guard layer maps the proposed command to an executable command:
\begin{equation}
    u_k = \gamma(u_k^p, o_k, G, u_{k-1}),
\end{equation}
where $\gamma$ enforces safety stops, command bounds, rate limits, enable
flags, and fallback behavior. The backend then applies the command:
\begin{equation}
    DT_{k+1} = A(DT_k, u_k, \Delta t),
\end{equation}
where $A$ is implemented by a simulator, a MATLAB/Simulink bridge, a live UDP
interface, or another runtime adapter. Thus, the policy selects $u_k^p$, while
the guard determines whether and how it becomes $u_k$.

\subsection{Component Interfaces}

The backend interface contains five operations:
\begin{equation}
\begin{split}
    &\texttt{reset()} \rightarrow o, \quad
    \texttt{observe()} \rightarrow o, \\
    &\texttt{apply\_command}(u), \quad
    \texttt{advance}(\Delta t) \rightarrow o, \quad
    \texttt{close()}.
\end{split}
\end{equation}
This interface is intentionally small. It is sufficient for mock simulation,
step-based simulator integration, and live communication. In live UDP mode,
the \texttt{advance} operation does not force the simulator to step; instead,
it waits until the live digital twin publishes an observation whose simulation
timestamp has advanced sufficiently.

The policy interface contains one operation:
\begin{equation}
    \texttt{choose\_action}(o, G) \rightarrow u^p.
\end{equation}
This allows a deterministic controller, optimization routine,
reinforcement-learning policy, or language-model policy to be substituted
without modifying the runner.

\section{Contracts and Guarded Execution}

\subsection{Observation, Command, and Goal Contracts}

The framework uses explicit contracts for observations, commands, and mission
goals. Table~\ref{tab:contracts} summarizes the case-study contract. The
observation is not intended to contain every simulator variable; it exposes
the variables required for supervisory decision-making and safety monitoring.
Additional internal variables remain inside the digital twin model.

\begin{table*}[t]
\caption{State-action-goal contract used in the battery-assisted dc bus case study.}
\label{tab:contracts}
\centering
\small
\begin{tabular}{|p{0.16\textwidth}|p{0.28\textwidth}|p{0.48\textwidth}|}
\hline
\textbf{Contract} & \textbf{Fields} & \textbf{Purpose} \\
\hline
Observation &
\texttt{sim\_time\_s}, \texttt{soc\_pct}, \texttt{v\_bus\_v},
\texttt{v\_batt\_v}, \texttt{i\_batt\_a}, \texttt{p\_batt\_w},
\texttt{p\_load\_w}, \texttt{p\_source\_w}, source flags, fault flags &
Expose the compact digital twin state needed for supervisory decisions,
constraint checks, logging, and backend-independent policy execution. \\
\hline
Command &
\texttt{mode}, \texttt{p\_batt\_cmd\_w}, \texttt{enable}, \texttt{reason} &
Restrict the autonomous supervisor to high-level commands:
\texttt{CHARGE}, \texttt{HOLD}, \texttt{DISCHARGE}, or \texttt{SAFE}.
Positive battery power discharges into the bus; negative power charges the
battery. \\
\hline
Goal &
SOC target band, SOC hard limits, bus nominal voltage, bus soft band, bus hard
band, success hold time, maximum mission time, optional current limit &
Separate desired operating regions from hard safety constraints and mission
termination criteria. \\
\hline
\end{tabular}
\end{table*}

For the current case study, the target SOC band is 40--80\%, the SOC hard
limits are 10--95\%, the nominal bus voltage is 400 V, the bus soft band is
$\pm$5\% of nominal, and the bus hard band is $\pm$20\% of nominal. These
values are configurable; they are included here to make the case study
concrete.

\subsection{Guarded Runner Logic}

The runner enforces safety and validity checks that are deliberately outside
the policy:
\begin{itemize}
    \item a startup grace period for tolerating simulator initialization
    transients;
    \item hard-constraint checking over SOC, bus voltage, optional current
    limits, and fault flags;
    \item soft-constraint reporting for target-band violations;
    \item command clipping against a maximum battery-power magnitude;
    \item command rate limiting between supervisor ticks;
    \item disabled-command behavior that resets the held power command to
    zero; and
    \item fallback behavior for AI policy failures.
\end{itemize}
These mechanisms make the policy a proposal generator rather than the final
execution authority.

\begin{figure}[t]
\centering
\fbox{\begin{minipage}{0.95\columnwidth}
\small
\textbf{Guarded supervisory loop}\\
1) Reset backend and observe digital twin.\\
2) Check hard constraints; issue \texttt{SAFE} on violation.\\
3) Record soft violations.\\
4) If goal is held for the required duration, return success.\\
5) During startup grace, hold. Otherwise, call policy.\\
6) Clip and rate-limit command.\\
7) Apply command, log event, and advance to next observation.\\
8) Stop on success, safety stop, timeout, backend failure, or strict policy failure.
\end{minipage}}
\caption{Compact algorithmic view of the guarded autonomous supervisory loop.}
\label{fig:algorithm}
\end{figure}

\subsection{Mission Outcomes}

The runner returns one of five statuses: \texttt{running},
\texttt{succeeded}, \texttt{timed\_out}, \texttt{safety\_stopped}, or
\texttt{failed}. This status set distinguishes control performance, safety
termination, infrastructure failure, and incomplete missions. Each decision
event stores the tick index, observation, command, active violations, and
status. This design supports later analysis without requiring the researcher
to reconstruct decisions from console output.

\section{Policy and Guard Layer}

\subsection{Deterministic Baseline Policy}

The baseline policy is a deterministic SOC and bus-voltage band policy. It
first checks active fault flags. If faults are active, it requests a safe
command. If bus voltage is below the soft lower limit and the battery is above
the hard reserve, it requests battery discharge to support the bus. If bus
voltage is above the soft upper limit and the battery is not near the upper
target, it requests charging to absorb power. If SOC is below the target band
and the source is available, it requests charging. If SOC is above the target
band, it requests discharge. Otherwise, it holds.

This policy is intentionally simple. Its role is not to represent a final
optimal controller. Its role is to provide a transparent baseline and fallback
policy for framework development.

\subsection{AI and LLM Policy Interface}

The same policy interface can be implemented by an AI policy. In the current
implementation, an OpenAI-compatible chat-completion client can be used as a
high-level supervisory policy. The model receives a structured JSON payload
containing the goal, observation, and required response schema. It must return
one compact JSON object:

\begin{figure}[t]
\centering
\fbox{\begin{minipage}{0.95\columnwidth}
\small
\texttt{\{}\\
\quad\texttt{"mode": "CHARGE | HOLD | DISCHARGE | SAFE",}\\
\quad\texttt{"p\_batt\_cmd\_w": 0.0,}\\
\quad\texttt{"reason": "short explanation"}\\
\texttt{\}}
\end{minipage}}
\caption{Required response schema for an LLM-based supervisory policy.}
\label{fig:llm_schema}
\end{figure}

The language model is not allowed to command duty ratio, switching states, or
direct actuator values. It only proposes the high-level supervisory command.
The runner parses the JSON, converts it to the command object, clips the
command, applies rate limits, and handles failures.

If the AI policy fails to produce valid JSON or if the model call fails, the
framework can fall back to the deterministic baseline policy. In strict
experimental mode, the same failure can instead terminate the run and record
the policy failure. This option separates control robustness experiments from
AI-output-validity experiments.

\subsection{Trace Logging}

For AI policies, the framework can log each model request, response, and
fallback event as JSON lines. This enables post-run analysis of AI behavior,
response validity, and policy rationale. Trace logging is separate from
mission logging so that control events and model-interaction events can be
analyzed independently.

\section{Backend and Communication Layer}

\subsection{Backend Modes}

The backend abstraction supports progressive development from offline testing
to live digital twin operation. Table~\ref{tab:backends} summarizes the
current backend modes.

\begin{table}[t]
\caption{Backend modes supported by the framework.}
\label{tab:backends}
\centering
\small
\begin{tabular}{|p{0.31\columnwidth}|p{0.57\columnwidth}|}
\hline
\textbf{Backend} & \textbf{Role} \\
\hline
Mock battery-bus simulator &
Deterministic approximation for runner, policy, guard, and report testing. \\
\hline
MATLAB/\newline Simulink bridge &
Step-based execution that initializes the model, applies command variables,
steps simulation time, and extracts observations. \\
\hline
Live UDP backend &
Interface to a continuously running Simulink model that publishes observations
and receives supervisory command packets. \\
\hline
\end{tabular}
\end{table}

The mock backend is not intended to replace the Simulink model. It exists so
the runner, policy interface, guard layer, command limiting, mission status
logic, and report generation can be tested before live model integration.

The MATLAB/Simulink backend uses MATLAB Engine to initialize the model, apply
command variables, step the simulation, and extract observations. The bridge
expects command variables such as \texttt{P\_batt\_cmd}, \texttt{enable\_cmd},
and \texttt{mode\_cmd}. Observations are returned as JSON to avoid fragile
MATLAB-to-Python struct conversion.

The live UDP backend supports a continuously running Simulink model. In this
mode, Simulink sends observation packets to the Python runner and receives
command packets from it. The runner does not call MATLAB Engine to advance the
simulation; instead, it waits for live observation packets and resends the
latest command at a configurable interval.

\subsection{UDP Protocol}

The UDP packet schema is compact JSON with a protocol name, version, packet
kind, sequence number, timestamp, and command or observation payload. Commands
include both string modes and numeric mode codes:
\begin{equation}
\begin{array}{ll}
\texttt{CHARGE}=-1, & \texttt{HOLD}=0,\\
\texttt{DISCHARGE}=1, & \texttt{SAFE}=99.
\end{array}
\end{equation}
The Simulink command receiver can include a stale-command timeout. If no fresh
command arrives within the allowed time, the model can disable the command
path and enter a safe mode. This creates a second safety layer inside the
digital twin runtime.

\section{Case Study: Battery-Assisted DC Bus}

\subsection{System Description}

The framework is demonstrated using a battery-assisted dc load-bus digital
twin. The digital twin includes a dc bus, source, load, battery, and
converter-level control path. The autonomous supervisor is responsible for
maintaining SOC within a target band while keeping bus voltage within a
desired range. The supervisor does not directly control the converter duty
ratio. Instead, it sends a battery-power command to the digital twin, and the
model's internal controller maps that command to converter behavior.

This case study is representative of a broader class of power-system digital
twin problems in which supervisory decisions must coordinate energy storage,
load support, source limits, and safety constraints. The same framework can
also be mapped to a laboratory dc microgrid with two parallel
converter-interfaced sources, droop-based load sharing, a common dc bus, and a
variable resistive load. In that setting, the observation and command
contracts would be extended to include source currents, droop coefficients,
load-sharing ratios, and bus-voltage regulation objectives.

\subsection{Implementation Summary}

The implementation includes the following artifacts:
\begin{itemize}
    \item \texttt{AutonomousRunner}, which implements the guarded supervisory
    loop;
    \item \texttt{DTObservation}, \texttt{DTCommand}, and \texttt{DTGoal},
    which define the state-action-goal contract;
    \item \texttt{SocBandPolicy}, which provides a deterministic baseline
    policy;
    \item \texttt{LLMControlPolicy}, which provides an OpenAI-compatible AI
    policy;
    \item \texttt{MockBatteryBusBackend}, \texttt{SimulinkBackend}, and
    \texttt{UdpLiveSimulinkBackend}, which support multiple runtime modes; and
    \item reporting utilities that generate CSV and Markdown audit logs.
\end{itemize}

The initial mission objective is intentionally modest: maintain SOC inside the
target band, maintain bus voltage inside the soft bus-voltage band, stop
safely if SOC or bus voltage crosses hard limits, and record each supervisory
decision and constraint violation. This objective validates the framework
before introducing more complex optimization objectives such as energy
management, adaptive droop coordination, or predictive scheduling.

\subsection{Mission Evidence}

The current prototype produces mission reports containing mission status,
mission reason, number of supervisor decisions, simulation time range, SOC
start and end values, bus-voltage start and end values, telemetry extrema,
commanded-power extrema, mode counts, constraint violations, and a decision
timeline. This format provides a direct audit trail for each run. For example,
nominal mock-backend runs show repeated \texttt{HOLD} decisions when SOC and
bus voltage are inside the target band. Hard-limit scenarios show deterministic
transition to \texttt{SAFE} mode. Live UDP trials demonstrate that the same
runner can interact with a continuously running digital twin through a packet
interface.

Because this is a framework paper, these demonstrations support feasibility,
modularity, guard behavior, and auditability rather than claiming universal
control optimality.

\section{Evaluation Plan}

The evaluation should answer four framework-level questions: (i) whether the
framework preserves safety constraints independently of the active policy;
(ii) whether different policies can be substituted through the same interface;
(iii) whether the same runner can operate over mock, step-based Simulink, and
live UDP backends; and (iv) whether the reporting layer produces enough
evidence to audit each decision.

\begin{table*}[t]
\caption{Planned scenario matrix for the final framework evaluation.}
\label{tab:scenarios}
\centering
\small
\begin{tabular}{|p{0.18\textwidth}|p{0.24\textwidth}|p{0.29\textwidth}|p{0.15\textwidth}|}
\hline
\textbf{Scenario} & \textbf{Initial condition or disturbance} &
\textbf{Expected framework behavior} & \textbf{Primary metric} \\
\hline
Nominal in-band operation & SOC and bus voltage inside target bands &
Hold command, no violations & Correct hold behavior \\
\hline
Low SOC with source available & SOC below target, source available &
Charge command subject to limits & Correct mode and bounded power \\
\hline
Low bus voltage & Bus voltage below soft band, SOC above reserve &
Discharge for bus support & Voltage-support response \\
\hline
Hard SOC violation & SOC at or beyond hard reserve or maximum &
Immediate safe stop & Safety-stop correctness \\
\hline
Communication interruption & Missing or stale UDP observations or commands &
Timeout, fallback, or safe behavior & Communication robustness \\
\hline
AI malformed output & Invalid JSON or invalid mode from policy &
Fallback or strict failure & Policy-failure handling \\
\hline
Backend substitution & Same scenario across mock and Simulink backends &
Same runner and policy interface & Backend portability \\
\hline
\end{tabular}
\end{table*}

Metrics should include task success rate, hard-constraint violation count,
soft-constraint violation count and duration, safe-stop correctness, command
validity rate, command clipping frequency, command rate-limit activation
frequency, AI fallback rate, communication timeout count, decision latency,
audit-log completeness, and backend portability across scenarios. These
metrics are aligned with the framework claim: the architecture should capture,
bound, and report each run consistently across policies and backends.

Recommended comparisons include deterministic baseline policy with guard
layer, AI policy with guard layer, AI policy in strict mode, guarded policy
with command rate limiting disabled in offline simulation, mock backend versus
Simulink backend, and step-based Simulink bridge versus live UDP
communication. An unguarded AI policy should not be run against a live or
physical system; if included, it should be restricted to offline simulation to
illustrate the value of guard enforcement.

\section{Discussion}

The proposed framework addresses a gap between digital twin modeling and
autonomous control. Many digital twin architectures emphasize representation,
synchronization, and visualization. Many control architectures emphasize
plant-level control design. The proposed framework focuses on the layer
between them: the guarded autonomous supervisor that treats the digital twin
as an executable runtime while preserving explicit safety and audit contracts.

The framework is especially useful when experimenting with AI-based policies.
Language models and other learned policies may offer flexible reasoning,
natural-language rationales, or adaptation to structured context. However,
they should not be treated as safety-critical controllers. In this framework,
the AI policy is an interchangeable proposal generator. It can suggest a
high-level command, but the runner checks constraints, clips commands,
rate-limits power, and decides whether to continue or stop.

The backend abstraction also improves development workflow. Researchers can
begin with a deterministic mock backend to validate runner logic and scenario
definitions. They can then use MATLAB/Simulink engine integration for
controlled step-based simulation. Finally, they can move to live UDP
communication when the digital twin is running continuously. This progression
reduces the risk of debugging AI behavior, communication behavior, and
plant-model behavior simultaneously.

\section{Limitations and Future Work}

The present framework has several limitations. First, the case-study objective
is intentionally simple. Maintaining SOC and bus-voltage bands demonstrates
the architecture but does not fully exploit predictive digital twin
capabilities. Future work should include forecast-aware energy management,
adaptive load sharing, and multi-objective optimization. Second, the current
command space is compact. More complex microgrid scenarios may require
commands for droop coefficients, source dispatch, load shedding,
operating-mode transitions, or fault-isolation decisions. Third, the current
implementation focuses on supervisory control in simulation and live digital
twin environments. Physical hardware integration would require additional
safety certification, real-time guarantees, hardware interlocks, and operator
approval pathways. Fourth, AI-policy evaluation remains preliminary. Final
claims about AI control quality require a larger scenario set, controlled
baselines, repeated trials, and analysis of malformed outputs, latency, and
fallback frequency.

Future work will extend the framework in four directions: a richer microgrid
case study with adaptive droop and load-sharing control; scenario libraries
for repeatable autonomous digital twin benchmarking; formal policy-safety
contracts for learned and LLM-based supervisors; and integration with
operator-in-the-loop approval and explanation workflows.

\section{Conclusion}

This paper proposed a modular framework for guarded autonomous digital twin
supervisory control. The framework defines a structured observation, goal,
command, policy, backend, guard, and reporting architecture. It separates
high-level policy proposals from deterministic safety enforcement and backend
execution. It supports deterministic and AI-based policies through a common
interface, and it can operate over mock simulation, MATLAB/Simulink engine
control, and live UDP communication. A battery-assisted dc bus digital twin
case study demonstrates how the framework can be instantiated for power-system
supervisory control while preserving auditability and safety constraints.

The main contribution is not a single controller tuned for one scenario.
Rather, the contribution is a reusable architecture for developing, testing,
and evaluating autonomous digital twin controllers. By making the
state-action contract explicit, keeping safety outside the policy, and
recording each decision, the framework provides a practical path from passive
digital twins toward guarded autonomous digital twin operation.

\section*{Acknowledgment}

TODO: Add funding, laboratory, collaborator, and data/model acknowledgments
required by the target IEEE special call.

\begin{thebibliography}{00}

\bibitem{dt_foundation}
TODO: Add a verified foundational digital twin reference in IEEE format.

\bibitem{dt_survey}
TODO: Add a verified digital twin survey reference in IEEE format.

\bibitem{dt_power}
TODO: Add a verified digital twin reference for power systems or microgrids.

\bibitem{dt_control}
TODO: Add a verified autonomous or control-oriented digital twin reference.

\bibitem{guarded_autonomy}
TODO: Add a verified guarded autonomy, safety filter, or safe supervisory
control reference.

\bibitem{llm_agents}
TODO: Add a verified reference for LLM agents or tool-using policies if the
LLM policy remains in the final submission.

\bibitem{simulink}
TODO: Add a verified model-based design, MATLAB/Simulink, or co-simulation
reference if required by the venue.

\bibitem{microgrid_control}
TODO: Add a verified dc microgrid or battery energy storage supervisory
control reference.

\bibitem{local_hierarchical_dt}
TODO: Add prior hierarchical or image-based digital twin work from the local
project literature set.

\bibitem{local_dc_microgrid_dt}
TODO: Add prior dc microgrid digital twin forecasting or control work from the
local project literature set.

\end{thebibliography}

\end{document}
